\begin{textsample}{POS Dim 4 – ac – Score 5.00 – ac/ac\_ef\_23.txt}  \label{ex:f4_pos_246}
3.
PARTE DIVERSIFICADA E ESPECIFICIDADES DO \textbf{ESTADO} DO ACRE O presente documento norteador deve apresentar, além de conhecimentos articulados com as competências gerais e específicas do componente curricular, conhecimentos que abordem a parte singular ou a chamada parte diversificada, que englobe temáticas necessárias, respeitando características regionais e locais da sociedade acreana.
No caso do ensino da Geografia, é importante salientar a abordagem geográfica e política do Acre, a materialização das lutas e disputas no espaço geográfico, configurando o espaço socioeconômico na formação cartográfica no que hoje temos como nosso espaço territorial, mostrando que nenhum \textbf{estado} tem uma formação igual ; mas surge de diferentes modos e condições.
Assim, é importante abordar os embates coloniais, as relações de \textbf{fronteiras}, os tratados políticos e as questões relacionadas à expansão territorial.
Também, de igual modo, no sentido de relevância, é importante trazer o conhecimento do ciclo da borracha, a vinda dos nordestinos e o processo \textbf{migratório} como um todo, a agricultura, a hidrografia e a biodiversidade do \textbf{estado} do Acre.
E não menos fundamental, o estudo do povo \textbf{indígena} na região acreana e suas contribuições para formação e cultura do povo acreano.
Nesse sentido, a Geografia, na sua parte diversificada, está relacionada à necessidade de se conhecer o espaço geográfico acreano.
Este pode ser entendido como o espaço produzido pelo \textbf{homem} e que está em constante transformação, ao longo do tempo.
Podemos dizer, então, que o espaço geográfico possui um caráter histórico-político e, por isso, é importante seu estudo para compreensão das características da ação humana sobre o meio em que vive.
Sendo, portanto, campo de estudo desta parte toda a dinâmica da natureza e social da Terra.

% matched lemmas: estado, fronteira, homem, indígena, migratório
\end{textsample}
